% Hodina 25 – revize podle komentářů k úvodu, přehledové tabulce a S1.
% Příprava k hodině 25 – přepracovaná po hodinách 21–24.
% Rozpis: 1_rocnik/tematicke_plany/Rozpis_vyuky_MAT_1_rocnik.tex, hodina 25 (TP 2.1–2.4).
% Učebnice: ucebnice/01_Zakladni_poznatky/Mocniny_a_odmocniny/Mocniny a Odmocniny.tex,
%           oddíl „Mocniny s přirozeným mocnitelem“ a navazující důkazy pravidel.
% Realisticky.cz: 010601 Mocniny s přirozeným mocnitelem, 010602 Vzorce pro mocniny I;
%           inspirace gradací, samostatnou prací a odvozováním pravidel z definice.
% Nové příklady jsou autorské, nikoli převzaté z Realisticky.cz.
\documentclass[11pt,a4paper]{article}
% Spolecny styl pro ucitelske pripravy MAT (Work/MAT)
% Verze 1.1, 2026-10-23
\NeedsTeXFormat{LaTeX2e}
\ProvidesPackage{priprava}[2026/10/23 v1.1 Sablona ucitelskych priprav MAT]

% Zakladni typografie a cestina
\RequirePackage[T1]{fontenc}
\RequirePackage[utf8]{inputenc}
\RequirePackage[czech]{babel}
\RequirePackage{lmodern}
\RequirePackage{microtype}
\RequirePackage{amsmath,amssymb}
\RequirePackage{xcolor}
\RequirePackage{enumitem}
\RequirePackage[most]{tcolorbox}
\RequirePackage{tabularx,booktabs,array}
\RequirePackage{fancyhdr}
\RequirePackage{needspace}
\RequirePackage{geometry}
\RequirePackage[hidelinks]{hyperref}

\geometry{a4paper,top=17mm,bottom=18mm,left=23mm,right=23mm}
\definecolor{notegray}{gray}{0.39}
\setlength{\parindent}{0pt}
\setlength{\parskip}{2.5pt plus 1pt minus 0.5pt}
\setlist[itemize]{topsep=2pt,itemsep=1pt,parsep=0pt,leftmargin=1.45em}
\setlist[enumerate]{topsep=2pt,itemsep=2pt,parsep=0pt,leftmargin=1.65em}
\setcounter{secnumdepth}{0}
\renewcommand{\arraystretch}{1.12}

% Zahlavi: vlevo rocnik | globalni tema, vpravo cislo: nazev hodiny
% Zapati zustava zcela prazdne.
\pagestyle{fancy}
\fancyhf{}
\renewcommand{\headrulewidth}{0.35pt}
\renewcommand{\footrulewidth}{0pt}
\setlength{\headheight}{14pt}
\newcommand{\nastavPripravu}[4]{%
  \fancyhead[L]{\small #1\enspace\textbar\enspace #2}%
  \fancyhead[R]{\small #3: #4}%
}

% Minimalisticke cernobile boxy. Kazda definice zvlast.
\tcbset{commonlesson/.style={enhanced,breakable=false,colback=white,colframe=black,
  colbacktitle=white,coltitle=black,boxrule=0.45pt,
  arc=0.8mm,outer arc=0.8mm,left=3mm,right=3mm,top=1.6mm,bottom=1.6mm,
  toptitle=1mm,bottomtitle=0.7mm,fonttitle=\bfseries,
  attach boxed title to top left={xshift=2.5mm,yshift=-1.15mm},
  boxed title style={colback=white,colframe=white,boxrule=0pt,left=1mm,right=1mm}}}
\newtcolorbox{zadaniBox}[1][]{commonlesson,title={Zadání},#1}
\newtcolorbox{prikladBox}[1][]{commonlesson,title={Příklad},#1}
\newtcolorbox{definiceBox}[1][]{commonlesson,title={Definice},#1}

% Spolecny box pro zadani a jeho reseni, bez mezery mezi castmi.
% Pouziti:
%   \begin{prikladReseniBox}
%     Zadani ulohy.
%     \tcblower
%     \begin{align*} ... \end{align*}
%   \end{prikladReseniBox}
% Spodni cast ma jemny sedy podklad a delici linku; lze v ni sazet
% vice radku matematiky. Obe casti maji jen jeden spolecny ramecek.
\tcbset{solutionlesson/.style={
  commonlesson,
  enhanced jigsaw,
  breakable,
  colbacklower=black!3,
  segmentation style={solid,draw=black!25,line width=0.35pt},
  fontlower=\small\color{notegray},
  before lower={\noindent\textbf{Řešení:}\enspace},
  leftlower=3mm,rightlower=3mm
}}
\newtcolorbox{prikladReseniBox}[1][]{solutionlesson,title={Příklad},#1}
\newtcolorbox{zadaniReseniBox}[1][]{solutionlesson,title={Zadání},#1}

% Pomucky pro autora pripravy: sedive poznamky; \reseni pro starsi pripravy
\newcommand{\poznamka}[1]{\par\begingroup\small\color{notegray}\itshape
  \textnormal{\textbf{Poznámka pro učitele.}}\enspace #1\par\endgroup}
\newcommand{\reseni}[1]{\par\begingroup\small\color{notegray}
  \textbf{Očekávané řešení:}\enspace #1\par\endgroup}
\newcommand{\blok}[1]{\needspace{5\baselineskip}\subsection*{#1}}
\newcommand{\vyrok}[1]{\uv{#1}}

\usepackage{tikz}
\usetikzlibrary{calc}
\nastavPripravu{1. ročník}{Mocniny a odmocniny}{25}{Mocnina s přirozeným exponentem}
\begin{document}

\poznamka{Jádro 45 minut: motivace a geometrie 5 min, definice a obrázek 7 min,
tabulka a zdůvodnění znamének 7 min, P1 7 min, P2 6 min,
S1 a společná kontrola 10 min, výstupní otázka 3 min. P3, S2 a R1 jsou rozšíření, nikoli povinný přídavek.}

\blok{Od geometrie k mocnině}
V hodině 22 jsme u Pythagorovy věty počítali obsahy čtverců.
Čtverec o straně $3\,\mathrm{cm}$ má obsah $3\cdot3=3^2=9\,\mathrm{cm}^2$.
Krychle o hraně $2\,\mathrm{cm}$ má objem $2\cdot2\cdot2=2^3=8\,\mathrm{cm}^3$.
Stejně jako \(4+4+4=3\cdot4\) zapisujeme opakované sčítání jako násobení,
\textbf{opakované násobení zkracujeme mocninou}: \(4\cdot4\cdot4=4^3\).
\textbf{Pořadí čísla a počtu opakování se liší:} v součinu \(3\cdot4\)
jsou tři sčítané čtyřky, zatímco v mocnině \(4^3\) jsou tři násobené čtyřky.

\begin{center}
\begin{tikzpicture}[x=.56cm,y=.56cm]
\draw[very thick] (0,0) rectangle (3,3);
\foreach \i in {1,2}{\draw[thin] (\i,0)--(\i,3);\draw[thin] (0,\i)--(3,\i);}
\draw[<->] (0,-.35)--node[below]{$3\,\mathrm{cm}$}(3,-.35);
\draw[<->] (-.35,0)--node[left]{$3\,\mathrm{cm}$}(-.35,3);
\node[right] at (4.0,1.5) {$S=3^2=9\,\mathrm{cm}^2$};
\end{tikzpicture}
\end{center}
\textbf{Pozor:} Mocnina není opakované sčítání. Zápis $3^4$ znamená
$3\cdot3\cdot3\cdot3=81$, nikoli $3\cdot4=12$.

\begin{definiceBox}
Nechť $a\in\mathbb R$ a $n\in\mathbb N=\{1,2,3,\ldots\}$.
\textbf{Mocnina s přirozeným exponentem} je součin $n$ stejných činitelů:
\[
 a^n=\underbrace{a\cdot a\cdot\ldots\cdot a}_{n\ \text{činitelů}}.
\]
Číslo $a$ nazýváme \textbf{základ mocniny (mocněnec)}, číslo $n$
\textbf{exponent (mocnitel)}; výraz $a^n$ je \textbf{mocnina}.
\end{definiceBox}

\begin{center}
\begin{tikzpicture}[>=latex,baseline=(current bounding box.center)]
\node[font=\Large] (mocnenec) at (0,0) {$a$};
\node[font=\normalsize] (mocnitel) at (.33,.43) {$n$};
\node[anchor=east,font=\small] (z) at (-1.5,-.48) {základ (mocněnec)};
\node[anchor=west,font=\small] (e) at (1.34,.68) {exponent (mocnitel)};
\draw[->,thick] (z.east)--(mocnenec.south west);
\draw[->,thick] (e.west)--(mocnitel.east);
\end{tikzpicture}
\end{center}
\poznamka{V obrázku oba názvy jen ukazujeme; nejde o dvě nové definice.
Zdůrazni, že exponent udává \emph{počet činitelů}, nikoli násobitele základního čísla.}

\blok{Pravidla pro mocniny — postupně doplňovaná tabulka}
Při sudém počtu záporných činitelů se minusy spárují, při lichém
jedno zbude. Dnes jen $n\ge1$; \textbf{tabulku postupně rozšíříme}
i o další exponenty.
\begin{center}
\small
\begin{tabularx}{\linewidth}{@{}>{\raggedright\arraybackslash}p{23mm}>{\raggedright\arraybackslash}p{55mm}X@{}}
\toprule
\textbf{Pravidlo} & \textbf{Zápis} & \textbf{Podmínky / význam} \\
\midrule
První mocnina & $a^1=a$ & $a\in\mathbb R$ \\
Nula & $0^n=0$ & $n\in\mathbb N$ \\
Jednička & $1^n=1$ & $n\in\mathbb N$ \\
Kladný základ & $a^n>0$ & $a>0$, $n\in\mathbb N$ \\
Záporný základ & $a^n>0$ & $a<0$, $n\in\mathbb N$ sudé \\
Záporný základ & $a^n<0$ & $a<0$, $n\in\mathbb N$ liché \\
Základ $-1$ & $(-1)^n=\begin{cases}1&n\text{ sudé},\\[-1mm]-1&n\text{ liché}.\end{cases}$ & $n\in\mathbb N$ \\
\bottomrule
\end{tabularx}
\normalsize
\end{center}
\begin{prikladReseniBox}[title={P1 — Stejná čísla, jiné závorky},breakable=false]
Porovnej $(-3)^2$, $-3^2$, $(-3)^3$ a $-(-3)^2$.
U každého výrazu nejprve nahlas urči, co je základem mocniny.
\tcblower
\begin{align*}
 (-3)^2&=(-3)\cdot(-3)=9, & -3^2&=-(3\cdot3)=-9,\\
 (-3)^3&=(-3)\cdot(-3)\cdot(-3)=-27,& -(-3)^2&=-9.
\end{align*}
Bez závorek se u $-3^2$ umocní pouze číslo $3$, vnější minus
zůstává. Zápisy $(-3)^2$ a $-3^2$ \emph{nejsou stejné}.
\end{prikladReseniBox}
\poznamka{Nechej porovnat také $(-1)^{100}$ a $(-1)^{101}$.
Student má zdůvodnit výsledek párováním záporných činitelů,
ne pouze pravidlem naučeným zpaměti.}

\blok{Umocňujeme i zlomky a celé výrazy v závorkách}
\begin{prikladReseniBox}[title={P2 — Co je celým základem?},breakable=false]
Vypočti $\left(-\frac23\right)^3$, $\left(\frac12\right)^4$
a $(-2+5)^3$. Výsledky nech ve zlomcích, pokud je to možné.
\tcblower
\begin{align*}
 \left(-\frac23\right)^3
 &=\left(-\frac23\right)\left(-\frac23\right)\left(-\frac23\right)
 =-\frac8{27},\\
 \left(\frac12\right)^4
 &=\frac12\cdot\frac12\cdot\frac12\cdot\frac12=\frac1{16},\\
 (-2+5)^3&=3^3=27.
\end{align*}
Závorku nejprve vyhodnotíme. Například $(-2+5)^3$ \emph{není}
$(-2)^3+5^3$; na konkrétních číslech dostaneme $27\ne117$.
\end{prikladReseniBox}

\blok{Samostatná práce — povinné jádro}
\begin{zadaniReseniBox}[title={S1 — Spočti mocniny},breakable=false]
\textbf{Spočti mocniny.} Všímej si závorek a toho, co přesně umocňuješ.
\begin{center}
\renewcommand{\arraystretch}{1.85}
\begin{tabularx}{\linewidth}{@{}XXXXX@{}}
a) $6^1$ & b) $(-2)^4$ & c) $-2^2$ & d) $(-2)^5$ & e) $-(-3^2)$ \\
f) $\left(\frac1{10}\right)^3$ & g) $0^{1415}$ & h) $1^{2011}$ & i) $-1^{1918}$ & j) $(-1)^{1620}$ \\
\end{tabularx}
\end{center}
\tcblower
\begin{enumerate}[label=\alph*),itemsep=.5pt]
\item $6^1=6$.
\item $(-2)^4=16$ (sudý exponent).
\item $-2^2=-(2^2)=-4$ (minus je před mocninou).
\item $(-2)^5=-32$ (lichý exponent).
\item $-(-3^2)=-(-9)=9$ (vnější minus změní znaménko).
\item $\left(\frac1{10}\right)^3=\frac1{1000}$.
\item $0^{1415}=0$.
\item $1^{2011}=1$.
\item $-1^{1918}=-(1^{1918})=-1$ (bez závorky kolem $-1$).
\item $(-1)^{1620}=1$ (sudý exponent záporného základu).
\end{enumerate}
\end{zadaniReseniBox}
\poznamka{Při opravě vyber dvě chybná řešení bez jmen a polož otázku
„Co je základ?“ Nikdy neopravuj jen znaménko bez vysvětlení závorek.}

\blok{Výstupní otázka na závěr}
\begin{zadaniReseniBox}[title={E1 — Bez kalkulačky, 2 minuty},breakable=false]
Rozhodni a zdůvodni: $(-5)^6=5^6$ a $(-5)^7=5^7$.
Které tvrzení platí a proč?
\tcblower
První platí (sudý počet záporných činitelů), druhé ne:
$(-5)^7=-5^7<0$, zatímco $5^7>0$.
\end{zadaniReseniBox}

\blok{Dobrovolné rozšíření a rezerva}
\begin{zadaniReseniBox}[title={R1 — Odhad, závorky a protipříklad},breakable=false]
\begin{enumerate}[label=\alph*)]
 \item Co je větší: $2^{10}$, nebo $10^2$? Vysvětli, proč neplatí obecné pravidlo
 $a^n=an$.
 \item Porovnej $\left(\frac13\right)^2$ a $\left(\frac13\right)^3$.
 Proč při další kladné mocnině zlomku menšího než $1$ hodnota klesá?
 \item Rozhodni, zda vždy platí $(a+b)^2=a^2+b^2$.
 Jestli ne, najdi protipříklad.
 \item Zapiš \emph{pouze pomocí násobení} číslo $(-2)^3\cdot(-2)^2$.
 Kolik záporných činitelů celkem dostaneš? Vzorec pro součin mocnin zatím nepoužívej.
\end{enumerate}
\tcblower
\begin{enumerate}[label=\alph*)]
 \item $2^{10}=1024>100=10^2$. Mocnina opakuje násobení, nikoli běžný
     součin základu a exponentu.
 \item $\frac19>\frac1{27}$. Násobení kladným číslem menším než $1$ hodnotu zmenší.
 \item Ne: pro $a=b=1$ je vlevo $(1+1)^2=4$, vpravo $1^2+1^2=2$.
 \item Pět záporných činitelů:
     $(-2)(-2)(-2)\cdot(-2)(-2)$; součin je záporný, $-32$.
     Zápis připravuje myšlenku hodiny 26.
\end{enumerate}
\end{zadaniReseniBox}
\poznamka{Rezerva 10–15 minut pro rychlou skupinu. U čtvrté úlohy
neprozrazuj následující vzorec; nech studenty, aby jej sami objevili.}

\blok{Domácí procvičení}
\begin{zadaniReseniBox}[title={D1 — Přibližně 15 minut},breakable=false]
\begin{enumerate}[label=\alph*)]
 \item Vypočti $(-5)^2$, $-5^2$ a $-(-5)^3$.
 \item Vypočti $\left(-\frac23\right)^3$ a $\left(\frac14\right)^2$.
 \item Urči $0^{15}$, $1^{25}$ a $(-1)^{25}$.
 \item Vypočti $3^2-2\cdot(-2)^3$.
 \item Porovnej $(-4)^3$ a $-4^3$; $\left(\frac13\right)^2$
    a $\left(\frac13\right)^3$.
 \item Najdi chybu: $(-2+3)^2=(-2)^2+3^2$. Spočti obě strany.
\end{enumerate}
\tcblower
\begin{enumerate}[label=\alph*)]
 \item $25$, $-25$, $125$.
 \item $-\frac8{27}$ a $\frac1{16}$.
 \item $0$, $1$, $-1$.
 \item $9-2\cdot(-8)=25$.
 \item První dvojice je rovna ($-64=-64$); ve druhé $\frac19>\frac1{27}$.
 \item Levá strana $1$, pravá $4+9=13$; mocnina součtu není
     obecně součtem mocnin sčítanců.
\end{enumerate}
\end{zadaniReseniBox}

\end{document}
